\documentclass[10pt,a4paper]{amsart}
\usepackage{amsmath,amssymb,amsthm,graphicx}
% For commutative diagrams we use Xy-pic
\usepackage[all]{xy}

% We use 2cell for 2-commutative diagrams.
\xyoption{2cell}
% Theorem environments.
%
\theoremstyle{plain}
\newtheorem{theorem}[subsection]{Theorem}
\newtheorem{proposition}[subsection]{Proposition}
\newtheorem{lemma}[subsection]{Lemma}

\theoremstyle{definition}
\newtheorem{definition}[subsection]{Definition}
\newtheorem{example}[subsection]{Example}
\newtheorem{exercise}[subsection]{Exercise}
\newtheorem{situation}[subsection]{Situation}

\theoremstyle{remark}
\newtheorem{remark}[subsection]{Remark}
\newtheorem{remarks}[subsection]{Remarks}

\numberwithin{equation}{subsection}

% Macros
%
\def\lim{\mathop{\mathrm{lim}}\nolimits}
\def\colim{\mathop{\mathrm{colim}}\nolimits}
\def\Spec{\mathop{\mathrm{Spec}}}
\def\Hom{\mathop{\mathrm{Hom}}\nolimits}
\def\Ext{\mathop{\mathrm{Ext}}\nolimits}
\def\SheafHom{\mathop{\mathcal{H}\!\mathit{om}}\nolimits}
\def\SheafExt{\mathop{\mathcal{E}\!\mathit{xt}}\nolimits}
\def\Sch{\mathit{Sch}}
\def\Mor{\mathop{\mathrm{Mor}}\nolimits}
\def\Ob{\mathop{\mathrm{Ob}}\nolimits}
\def\Sh{\mathop{\mathit{Sh}}\nolimits}
\def\NL{\mathop{N\!L}\nolimits}
\def\CH{\mathop{\mathrm{CH}}\nolimits}
\def\proetale{{pro\text{-}\acute{e}tale}}
\def\etale{{\acute{e}tale}}
\def\QCoh{\mathit{QCoh}}
\def\Ker{\mathop{\mathrm{Ker}}}
\def\Im{\mathop{\mathrm{Im}}}
\def\Coker{\mathop{\mathrm{Coker}}}
\def\Coim{\mathop{\mathrm{Coim}}}

% Boxtimes
%
\DeclareMathSymbol{\boxtimes}{\mathbin}{AMSa}{"02}

%
% Macros for moduli stacks/spaces
%
\def\QCohstack{\mathcal{QC}\!\mathit{oh}}
\def\Cohstack{\mathcal{C}\!\mathit{oh}}
\def\Spacesstack{\mathcal{S}\!\mathit{paces}}
\def\Quotfunctor{\mathrm{Quot}}
\def\Hilbfunctor{\mathrm{Hilb}}
\def\Curvesstack{\mathcal{C}\!\mathit{urves}}
\def\Polarizedstack{\mathcal{P}\!\mathit{olarized}}
\def\Complexesstack{\mathcal{C}\!\mathit{omplexes}}
% \Pic is the operator that assigns to X its picard group, usage \Pic(X)
% \Picardstack_{X/B} denotes the Picard stack of X over B
% \Picardfunctor_{X/B} denotes the Picard functor of X over B
\def\Pic{\mathop{\mathrm{Pic}}\nolimits}
\def\Picardstack{\mathcal{P}\!\mathit{ic}}
\def\Picardfunctor{\mathrm{Pic}}
\def\Deformationcategory{\mathcal{D}\!\mathit{ef}}
\usepackage{cleveref}
\crefname{originalreference}{original reference}{original references}
\Crefname{originalreference}{Original reference}{Original references}
\makeatletter
\newlabel{lemma-one-constructible}{{095Z}{source}{}{sourcepointer-1}{}}
\newlabel{lemma-one-constructible@cref}{{[originalreference][0][]095Z}{[1][0][]source}{}{}{}}
\makeatother
\title{Bounded above complexes with constructible cohomology are represented by bounded above complexes of finite sums of extension by zero of the constant sheaf on affine objects}
\author{The Stacks Project Authors}
\date{stacks-project:etale-cohomology.tex:local-snapshot-20261011; GNU FDL 1.2 or later}
\begin{document}
\maketitle
\noindent\textit{Editorial note.} All delivered mathematical text is the author's own, copied from the pinned Stacks Project chapter file etale-cohomology.tex; no statement or proof step was rewritten, reordered or removed.\par\medskip
\noindent\textit{Editorial note.} The counted proof is the author's complete argument: the two immediate implications, the downward induction on the cohomological degree with the construction of the truncated representing complexes, and the final gluing step via the octahedral axiom.\par\medskip
\noindent\textit{Editorial note.} The label lemma-bounded-above-c names the counted statement, and the proof environment immediately following it in the source is delivered verbatim as its proof.\par\medskip
\noindent\textit{Editorial note.} The counted statement is carried by the author's own lemma environment in the chapter source.\par\medskip
\noindent\textit{Editorial note.} Pre-open gate: python operations/agent-authoring/tools/precard\_gate.py --source etale-cohomology.tex --label lemma-bounded-above-c --start 13643 --end 13685 --local-id 3552 returned RESULT=OK (L3552 free, occupancy FREE, public interval clear). Candidate taken from lane-s4/halfB-priority.json.\par\medskip
\noindent\textit{Editorial note.} A same-named label exists in sibling Stacks chapters that restate this material over other sites; no task card of this chapter counts this label, and the counted interval overlaps no public or card interval of the chapter.\par\medskip
\noindent\textit{Editorial note.} The statement and proof are copied from the pinned chapter file; the cited inputs are kept as external pointers with their canonical Stacks tags.\par\medskip
\noindent\textit{Allowed standard prerequisites (editorial).}\begin{itemize}
\item One constructible cohomology sheaf is generated by finitely many affine extensions by zero. Same chapter, the lemma producing a finite direct sum of j\_\{U!\} of the constant sheaf with U affine mapping onto a given constructible sheaf of Lambda-modules. Supplies, at each stage of the induction, the bounded-above piece F\textasciicircum{}a together with the map to the complex inducing the required isomorphism and surjection on cohomology.
\item The octahedral axiom TR4 of triangulated categories. Derived Categories, the axiom TR4 for triangulated categories. Supplies the compatibility of the shift-and-compose construction that turns the maps F\textasciicircum{}\{a-1\}[-a+1] -> L and L -> F\textasciicircum{}bullet into a map of complexes F\textasciicircum{}\{a-1\} -> ... -> F\textasciicircum{}b over K.
\end{itemize}\medskip
\ifdefined\hypertarget\hypertarget{sourcepointer-1}{}\fi
\noindent\textit{Source reference.} Original reference 095Z: sources/stacks-etale-cohomology/pinned-original-source.tex, source lines 13611--13611. Prerequisite pointer invoked at proof line(s) 13668; it names a standard fact of the same chapter that the counted argument uses, and every step of the counted proof built on it is retained verbatim in the delivered body.\par\medskip
\begin{lemma}
\label{lemma-bounded-above-c}
Let $X$ be a quasi-compact and quasi-separated scheme.
Let $\Lambda$ be a Noetherian ring. Let $K \in D^-(X_\etale, \Lambda)$. Then
the following are equivalent
\begin{enumerate}
\item $K$ is in $D_c(X_\etale, \Lambda)$,
\item $K$ can be represented by a bounded above complex
whose terms are finite direct sums of $j_{U!}\underline{\Lambda}$
with $U \in \Ob(X_\etale)$ affine,
\item $K$ can be represented by a bounded above complex
of flat constructible sheaves of $\Lambda$-modules.
\end{enumerate}
\end{lemma}

\begin{proof}
It is clear that (2) implies (3) and that (3) implies (1).
Assume $K$ is in $D_c^-(X_\etale, \Lambda)$.
Say $H^i(K) = 0$ for $i > b$. By induction on $a$
we will construct a complex $\mathcal{F}^a \to \ldots \to \mathcal{F}^b$
such that each $\mathcal{F}^i$ is a finite direct sum
of $j_{U!}\underline{\Lambda}$ with $U \in \Ob(X_\etale)$ affine
and a map $\mathcal{F}^\bullet \to K$ which induces an isomorphism
$H^i(\mathcal{F}^\bullet) \to H^i(K)$ for $i > a$ and a surjection
$H^a(\mathcal{F}^\bullet) \to H^a(K)$.
For $a = b$ this can be done by Lemma \ref{lemma-one-constructible}.
Given such a datum choose a distinguished triangle
$$
\mathcal{F}^\bullet \to K \to L \to \mathcal{F}^\bullet[1]
$$
Then we see that $H^i(L) = 0$ for $i \geq a$. Choose
$\mathcal{F}^{a - 1}[-a +1] \to L$ as in
Lemma \ref{lemma-one-constructible}. The composition
$\mathcal{F}^{a - 1}[-a +1] \to L \to \mathcal{F}^\bullet$
corresponds to a map $\mathcal{F}^{a - 1} \to \mathcal{F}^a$
such that the composition with $\mathcal{F}^a \to \mathcal{F}^{a + 1}$
is zero. By TR4 we obtain a map
$$
(\mathcal{F}^{a - 1} \to \ldots \to \mathcal{F}^b) \to K
$$
in $D(X_\etale, \Lambda)$. This finishes the induction step and the
proof of the lemma.
\end{proof}

\end{document}
